\section{Native residuals and trained direction controls}
\label{sec:2}
\subsection{A common construction, with backbone-specific operators}\label{a-common-construction-with-backbone-specific-operators}

Let \(k\) denote an injection site and its invocation, including recycle state.
At this invocation, let \(a_k,b_k\) be native factors, \(U_{k,0}\) the unmodified
update, and \(G_k\) the frozen factor-combination map used to construct a residual.
A sequence-conditioned adapter produces residue increments \((\delta a,\delta b)\):

\[
\Delta U_k=G_k(a_k+\delta a,b_k+\delta b)-G_k(a_k,b_k),
\qquad
U_k^{(R)}=U_{k,0}+\mathcal R_k\Delta U_k.
\tag{1}
\]

\(\mathcal R_k\) applies a fixed orthogonal matrix \(R_k\) to each pair's output
channels; the native arm uses \(R_k=I\). The chosen matrix is reused across
invocations at the selected site. Only the final increment projection is
zero-initialized, making every arm exactly equal to its backbone's baseline
function at initialization. Upstream adapter layers are normally initialized.
Backbone, decoder and PLM parameters remain frozen.

\(U_{k,0}\) means the unmodified update \textbf{at the current state}, not necessarily
one tensor cached for the whole forward. In AF2 and AtlasFold, live factors
can change with recycle state, including the effects of earlier adapter calls.
The Protenix experiment uses its cached query-only factors. These are different
instantiations of Eq. (1); the shared principle does not require a shared
normalization constant or structure loss.

For Protenix, the implemented residual decoder is
\(G_D(a,b)=[D W\operatorname{vec}(a\otimes b)+b_{\rm out}]/(D+0.001)\),
with \(D=508.5\), the training-set median MSA depth. The baseline \(U_q\) retains
its exact depth-one affine and normalization. \(D\) is a scalar in the residual
construction, not homolog information supplied at inference. For OpenFold,
\(G_k\) uses the selected layer's native masked OPM and pair-specific denominator;
the Protenix scalar is not transferred.

For the non-OPM AtlasFold extension \citep{atlasfold_code}, the frozen map combines projected factors
as \(G_k(a_i,b_j)=W_k[a_i-b_j;\,a_i\odot b_j]+b_k\).
Its native AtlasLM remains present. This extension has a fixed experimental
slot below, but its incomplete adapter results are not evidence in this draft.

\begin{table}[t]
\centering
\caption{Backbone-specific implementations of Eq.~(1). N/R: native and rotated factor heads. All adapters additionally receive frozen ESM2-35M layer-12 features (480 channels). AtlasFold adaptation results remain pending.}
\label{tab:interfaces}
\small
\begin{tabularx}{\linewidth}{@{}lXXX@{}}
\toprule
& Protenix-Mini / Tiny & AF2 / OpenFold & AtlasFold \\
\midrule
Site & Single MSA OPM boundary; pair stack retained & First main Evoformer OPM; all 48 blocks retained & First LM-stack product/difference update \\
Anchor & Cached query factors; $r=32$ & Live MSA factors; $r=32$ & Live sequence factors; $r=128$ \\
Normalization & Residual $D/(D+.001)$; exact depth-one baseline & Masked native row count plus layer epsilon & Native affine; no OPM depth constant \\
Native PLM & None in adapted default checkpoint & None & AtlasLM-3B retained \\
Parameters & N/R: 373,824; Generic: 378,272; G+: 378,144 & N/R: 373,824; G+: 378,144 & N/R: 423,168; G+: 415,008 \\
Loss & $4L_{\rm MSE}+4L_{\rm bond}$\newline $+4L_{\rm smooth}+.03L_{\rm dist}$ & $L_{\rm FAPE}+.3L_{\rm dist}$\newline $+L_{\rm torsion}$ & $.4L_{\rm dist}$\newline $+2(L_{\rm wMSE}+L_{\rm smooth})$ \\
Inference & 4 recycles; 5 diffusion steps; 1 sample; seed 101 & 4 trunk passes; no early stop; seed 20260921 & 5 trunk passes; 20/30 diffusion steps by length; 1 sample; seed 1 \\
\bottomrule
\end{tabularx}
\end{table}


The Protenix bond term is zero for the recorded empty protein-only bond masks.
Confidence losses are disabled; OpenFold also excludes masked-MSA loss.
AtlasFold retains native MLM conditioning with probability 0.15. Its pinned
implementation executes \nolinkurl{num_recycles+1} trunk passes. Losses retain native
component definitions and supervision masks;
experimental structures enter labels and scoring, not sequence-only inputs.
The appendix specifies precision, optimizer settings and invocation semantics.

\subsection{What the intervention controls}\label{what-the-intervention-controls}

At a common factor input and writer parameter state,

\[
\|\mathcal R\Delta U\|_F=\|\Delta U\|_F,
\quad \sigma_j(RW)=\sigma_j(W),
\quad J_R^\top J_R=J^\top J.
\tag{2}
\]

Here \(J\) is the local residual writer Jacobian, holding its incoming state
fixed. The downstream responses \(J_{\rm down}J\) and
\(J_{\rm down}\mathcal R J\) can differ. Independently trained trajectories,
state-dependent anchors and the full recurrent network are not asserted to
have identical Jacobian spectra or optimization geometry. The fixed baseline
and downstream weights are never rotated.

For OPM, the full residual contains a first-order term
\(T=\mathcal W(\delta a\otimes b+a\otimes\delta b)\) and quadratic increment
term \(Q=\mathcal W(\delta a\otimes\delta b)\). The first-order arm retains
\(T\) only. In the Protenix setting without an additional pair-dependent gate,
each output-channel residual has rank at most \(2r=64\); this bound need not
survive downstream nonlinear processing. The quadratic term is a combination
of adapter increments, not measured evolutionary covariance.

Generic predicts a pair residual using projected sequence features and clipped
relative sequence separation. G+ additionally receives the relevant native
factors. All heads share parameters across residue positions. G+ changes the
pair MLP width to approximately match parameter count; it is a competitive
explicit-access control, not a pure one-variable ablation. The N/R comparison
asks about orientation within a fixed family; Native-G+ asks which family
performs better. These contrasts need not have the same sign.
